% ═════════════════════════════════════════════════════════════
% Comparison: QR vs.\ L-BFGS
\label{sec:comparison}
% ═════════════════════════════════════════════════════════════

\noindent This chapter compares L-BFGS and structure-based Householder QR
under a common experimental protocol. Both solve the same regularized
least-squares problem, but their computational costs and sensitivity to the
spectrum differ. We therefore compare complete solve times together with
coefficient errors, optimality residuals, and termination status. All results
and both figures below come from the executed notebook
\texttt{04\_comparison.ipynb}.

\section{Experimental Protocol and Accuracy Measures}
\label{sec:comparison_protocol}

\noindent The initial problem is the standardized ML-CUP25 matrix with
$m=500$, $n=12$, and $\lambda=0.5$, using the same synthetic response vector
as the preceding experiments (seed 42). Consequently, accuracy here concerns
the numerical solution of the optimization problem, rather than prediction
on unseen data. L-BFGS uses the configuration selected in
Chapter~\ref{sec:experiments}: memory $\bar m=12$, BB1 scaling, exact line
search, no curvature-based restart, and $w_0=0$. Unless explicitly varied,
the relative gradient tolerance is $\varepsilon=10^{-14}$ and the iteration
budget is 2000. The curvature-pair acceptance safeguard remains active.

\noindent Each complete solver call is timed externally after one discarded
warm-up, with ten repetitions and a reproducibly shuffled method order.
The reported statistic is the median; the interval between the first and
third quartiles describes timing dispersion, not a confidence interval.
Initialization, factorization, and back substitution are included where
applicable. Data construction, SVD references, and accuracy diagnostics are
excluded. The run uses double precision, Python 3.13.9, NumPy 2.1.3, and one
BLAS thread on an ARM64 machine. These timings describe the present
implementations and should not be directly substituted for the internal
timings in the individual algorithm chapters.

\noindent The independent thin-SVD solution of
Eq.~\eqref{eq:qr_svd_reference} supplies a common numerical reference
$w_{\mathrm{SVD}}$. We use the forward error of
Eq.~\eqref{eq:qr_forward_error} and the relative gradient
\begin{equation}
\label{eq:comparison_gradient}
    e_w=\frac{\norm{w-w_{\mathrm{SVD}}}_2}{\norm{w_{\mathrm{SVD}}}_2},
    \qquad
    g_{\mathrm{rel}}=\frac{\norm{X(X^Tw-y)+\lambda^2w}_2}{\norm{Xy}_2}.
\end{equation}
Since $w_0=0$, the latter agrees with the stopping normalization in
Eq.~\eqref{eq:stop_rel}. It differs from $\eta_{\mathrm{opt}}$ in
Eq.~\eqref{eq:qr_scaled_optimality}, whose denominator contains the augmented
residual. A QR result is labelled as a direct solve; it is assessed with the
same metrics, without assigning it an iterative stopping tolerance.

\noindent The notebook also records fitted-vector discrepancies, data and
augmented residuals, and a scaled stationarity residual. Agreement in the
objective alone is insufficient near convergence. In particular, writing
$e=w-w_{\mathrm{SVD}}$ and $H=XX^T+\lambda^2I_m$ gives
\begin{equation}
\label{eq:comparison_energy}
    f(w)-f(w_{\mathrm{SVD}})
    =\nabla f(w_{\mathrm{SVD}})^Te+\tfrac12e^THe.
\end{equation}
The nonnegative quantity $e^THe/(2f(w_{\mathrm{SVD}}))$ avoids subtracting
nearly equal objectives, but equals the relative objective gap only when
the reference is an exact minimizer. The SVD reference itself has
$g_{\mathrm{rel}}\approx2.0\times10^{-15}$ on the initial instance.

\section{Accuracy and Runtime on ML-CUP25}
\label{sec:comparison_baseline}

\noindent Table~\ref{tab:comparison_baseline} shows that both algorithms
reach coefficient errors of order $10^{-15}$. L-BFGS satisfies the requested
tolerance in 12 iterations and is approximately $21.7$ times faster in this
run. Differences between forward errors at this scale do not establish a
meaningful accuracy advantage over a floating-point SVD reference.

\begin{table}[H]
\centering
\caption{Complete solve times and achieved accuracy on ML-CUP25
($m=500$, $n=12$, $\lambda=0.5$). Times and quartiles are in milliseconds.}
\label{tab:comparison_baseline}
\small
\begin{tabular}{lrrrr}
    \toprule
    \textbf{Method} & \textbf{Median} & $[Q_{0.25},Q_{0.75}]$
        & $e_w$ & $g_{\mathrm{rel}}$ \\
    \midrule
    L-BFGS & 0.512 & $[0.499,0.522]$ & $2.03\times10^{-15}$ & $4.26\times10^{-17}$ \\
    Structure-based QR & 11.114 & $[11.035,11.340]$ & $2.64\times10^{-15}$ & $6.49\times10^{-16}$ \\
    \bottomrule
\end{tabular}
\end{table}

\noindent Both solutions give $f(w)\approx0.092632$, a relative data
residual of $0.029902$, and a relative augmented residual of $0.161038$.
These nonzero residuals are compatible with optimality of the regularized
problem. In particular, minimizing the data residual alone would ignore the
penalty on $w$.

\subsection{Time at a Common Accuracy Requirement}
\label{sec:comparison_quality}

\noindent Figure~\ref{fig:comparison_time_accuracy} varies only the L-BFGS
tolerance, from $10^{-2}$ to $10^{-14}$. The three tightest tolerances all
produce the same 12-iteration solution; their small timing differences are
measurement variability. QR contributes one complete direct solve, rather
than an iterative convergence trajectory.

\begin{figure}[H]
    \centering
    \includegraphics[width=\textwidth]{Chapters/images/comparison_time_accuracy.png}
    \caption{Measured time--accuracy trade-off on the initial problem.
    Labels indicate requested L-BFGS tolerances. All sampled L-BFGS runs
    satisfy their tolerance; the square denotes the direct QR result.}
    \label{fig:comparison_time_accuracy}
\end{figure}

\noindent To compare the cost of reaching a common quality level, impose
$e_w\leq10^{-6}$ and $g_{\mathrm{rel}}\leq10^{-8}$. Among the sampled
configurations, the fastest qualifying L-BFGS run requests
$\varepsilon=10^{-8}$ and takes 10 iterations: its median time is
$0.388$ ms, with $e_w=1.41\times10^{-7}$ and
$g_{\mathrm{rel}}=8.69\times10^{-9}$. QR also meets both targets, taking
$11.114$ ms, so the measured time ratio is approximately $28.7$.
This is a comparison at common acceptance thresholds, rather than identical
achieved errors or a claim of globally optimal parameter tuning.

\section{Dimensional Scaling and Storage}
\label{sec:comparison_scaling}

\subsection{Increasing the Number of Unknowns}

\noindent The $m$ sweep follows the nested row-resampling construction of
Section~\ref{sec:qr_scalability}: original rows are retained, additional
rows are sampled with a small perturbation, and prefixes give
$m\in\{100,250,500,1000,2000,4000\}$ at fixed $n=12$ and $\lambda=0.5$.
Unlike the prescribed-spectrum study in Section~\ref{sec:exp_scaling_nm},
conditioning can change with $m$.

\begin{figure}[H]
    \centering
    \includegraphics[width=\textwidth]{Chapters/images/comparison_scaling_m.png}
    \caption{Complete solve times, coefficient errors, and L-BFGS iteration
    counts for increasing $m$. Dashed complexity guides are normalized at
    the smallest size and are not fitted asymptotic laws.}
    \label{fig:comparison_scaling_m}
\end{figure}

\noindent L-BFGS takes 12 iterations at every sampled size and meets its
tolerance throughout. Its time rises from $0.431$ ms at $m=100$ to
$1.745$ ms at $m=4000$, whereas QR rises from $1.801$ to $237.894$ ms.
The corresponding QR/L-BFGS ratio grows from $4.2$ to $136.3$.
Both methods retain forward errors below $10^{-14}$ over this sweep.

\noindent The growing runtime separation is consistent with the cost
models of Sections~\ref{sec:exp_cost} and~\ref{sec:qr_flop_count}:
\begin{equation}
\label{eq:comparison_cost}
    T_{\mathrm{L\text{-}BFGS}}=O\!\left(k(mn+\bar m m)\right),
    \qquad
    T_{\mathrm{QR}}=O\!\left((n+1)m^2\right).
\end{equation}
The first bound depends on the observed iteration count $k$; linear growth
in $m$ requires that this count remain bounded. Fixed overheads and the
finite range of sizes prevent interpreting the measured curves as exact
power laws.

\noindent Storage also favors L-BFGS in this regime. Counting only the
principal double-precision arrays gives $8(2\bar m m)$ bytes for its pair
capacity and $8(m^2+m(n+1))$ bytes for QR's dense triangular-factor array
and compact reflectors. At $m=4000$ these are $0.732$ and $122.467$ MiB,
respectively. Inputs, temporary arrays, Python objects, and convergence
histories are excluded; these figures are array counts, not measured peak
process memory.

\subsection{Increasing the Number of Spectral Directions}

\noindent A separate synthetic sweep fixes $m=500$, $\lambda=0.5$, and
$\bar m=12$, while increasing $n$ from 4 to 96. Singular values are
geometrically spaced between 1 and 10, keeping
$\kappa_2(A)\approx20.0$ and
$\kappa_2(H|_{\operatorname{range}(X)})=80.2$ constant. This differs from
the spectrum and memory study in Section~\ref{sec:exp_scaling_nm}.

\noindent The L-BFGS iteration counts are $4,8,12,33,78,113$ for
$n=4,8,12,24,48,96$, respectively. All runs satisfy the requested tolerance,
but fixed memory becomes less effective as more spectral directions are
introduced. At $n=96$, L-BFGS takes $10.255$ ms and QR $27.517$ ms, reducing
the timing advantage to approximately $2.7$. Their forward errors are
$1.91\times10^{-14}$ and $3.62\times10^{-15}$, respectively.
Thus the favorable behavior at $n=12$ should not be extrapolated with a
constant iteration count to larger $n$.

\section{Conditioning and a Difficult Synthetic Spectrum}
\label{sec:comparison_conditioning}

\noindent On ML-CUP25, decreasing $\lambda$ from $0.5$ to $10^{-8}$
raises the full augmented condition number to approximately
$7.6\times10^9$. Nevertheless, L-BFGS still converges in 12 iterations
with $e_w\approx2.1\times10^{-15}$. Structure-based QR instead gives
$e_w\approx7.7\times10^{-8}$, despite
$g_{\mathrm{rel}}\approx5.6\times10^{-16}$. This illustrates why a tiny
optimality residual does not guarantee accurate coefficients.

\noindent The subspace structure explains why the full condition number
alone does not predict L-BFGS behavior here. Starting from zero, in exact
arithmetic its iterates remain in $\operatorname{range}(X)$, where
\begin{equation}
\label{eq:comparison_restricted_condition}
    \kappa_2\!\left(H|_{\operatorname{range}(X)}\right)
    =\frac{\sigma_{\max}(X)^2+\lambda^2}
           {\sigma_{\min}(X)^2+\lambda^2}
\end{equation}
for full-column-rank $X$. This ratio remains about $10^3$ as
$\lambda\to0$ for the present data, while the full Hessian has additional
eigenvalues $\lambda^2$. The observed 12-step convergence is not a general
finite-termination guarantee for L-BFGS: the changing initial scaling,
limited memory, curvature safeguards, and rounding distinguish it from the
exact-arithmetic result discussed in Section~\ref{sec:finite_termination}.

\noindent To test a less favorable restricted spectrum, we construct one
matrix with $m=160$, $n=24$, and singular values geometrically spaced from
$10^{-4}$ to 1. Table~\ref{tab:comparison_stress} reports the results for
three regularization values, keeping the L-BFGS settings unchanged.

\begin{table}[H]
\centering
\caption{Synthetic spectral stress test. Times are medians in milliseconds;
the asterisk marks exhaustion of the 2000-iteration budget. QR has no
iterative stopping criterion.}
\label{tab:comparison_stress}
\small
\begin{tabular}{llrrrr}
    \toprule
    $\lambda$ & \textbf{Method} & \textbf{Time} & $k$ & $e_w$ & $g_{\mathrm{rel}}$ \\
    \midrule
    $10^{-4}$ & L-BFGS & 111.577 & $2000^*$ & $3.18\times10^{-5}$ & $9.23\times10^{-8}$ \\
               & Structure-based QR & 3.354 & -- & $3.95\times10^{-13}$ & $8.96\times10^{-14}$ \\
    $10^{-2}$ & L-BFGS & 2.997 & 57 & $3.92\times10^{-13}$ & $6.91\times10^{-15}$ \\
               & Structure-based QR & 3.272 & -- & $1.70\times10^{-14}$ & $1.17\times10^{-15}$ \\
    $1$       & L-BFGS & 0.319 & 9 & $4.62\times10^{-15}$ & $5.17\times10^{-16}$ \\
               & Structure-based QR & 3.273 & -- & $4.56\times10^{-15}$ & $5.72\times10^{-16}$ \\
    \bottomrule
\end{tabular}
\end{table}

\noindent At $\lambda=10^{-4}$, the restricted Hessian condition number
is approximately $5\times10^7$. L-BFGS reaches its budget without satisfying
the tolerance, while QR is about 33 times faster and gives a forward error
almost eight orders of magnitude smaller. QR's relative gradient is still
above $10^{-14}$, so its advantage should be stated in terms of the achieved
metrics, rather than attributed to satisfying the L-BFGS stopping test.
At $\lambda=10^{-2}$, L-BFGS converges in 57 iterations and the times are
similar, with QR more accurate in the coefficients. At $\lambda=1$, both
reach errors of order $10^{-15}$ and L-BFGS is faster again. The weakest
regularization in this synthetic test is the only L-BFGS configuration in
the recorded experiments that misses its requested tolerance.

\section{Conclusions}
\label{sec:comparison_conclusions}

\begin{itemize}
    \item \textbf{For the tested ML-CUP25 problems}, L-BFGS combines
    coefficient accuracy near working precision with shorter runtimes and
    substantially smaller principal storage requirements.
    \item \textbf{As the spectral dimension increases}, fixed L-BFGS memory
    requires more iterations, reducing its runtime advantage even at
    controlled conditioning.
    \item \textbf{For the difficult weakly regularized synthetic problem},
    structure-based QR is faster and more accurate than the budget-limited
    L-BFGS run. On the real-data small-$\lambda$ sweep, the accuracy ranking
    reverses.
\end{itemize}
\noindent The preferred solver therefore depends on the spectrum,
regularization, required accuracy, and available storage. These experiments
support a conditional choice between the two implementations, rather than
a uniform ranking of direct and iterative methods.
